%==============================================================================
% Chapter 5 -- Future work
%==============================================================================
\chapter{Future work}
\label{ch:future-work}

% -----------------------------------------------------------------------------
% The "expanding funnel": start from a very concrete immediate next step and
% widen to a multi-year research vision, closing with a Gantt chart. Edit the
% pgfgantt block at the end to match your real milestones and dates.
% -----------------------------------------------------------------------------

Having shown that we can model a whole symmetric assembly end to end
(Chapters~\ref{ch:methods}--\ref{ch:results}), my research programme now
\emph{expands}: from a precise, near-term objective to a broad, multi-year agenda.
The immediate goal is sharply defined; each subsequent direction opens the problem
a little wider.

\section{From unconditional to conditional generation: symmetric motif scaffolding}
\label{sec:motif-scaffolding}

Design-CP as reported performs \emph{unconditional} generation: it produces a
symmetric particle, but not a particle built around a \emph{specific} antigen. The
natural next step is \emph{conditional} generation, holding one or more functional
motifs (the antigen, or an epitope we wish to display) fixed while the model
designs a symmetric scaffold that presents them. In the symmetric setting this is
\emph{symmetric motif scaffolding}: the standard motif-scaffolding problem
\parencite{watsonNovoDesignProtein2023,didiFrameworkConditionalDiffusion2024},
recently generalised to atom-level motifs without known sequence indices
\parencite{ahern_atom_2025}, composed with the
point-group machinery of Section~\ref{sec:rfd3-symmetry}, so that a motif placed in
one ASU is presented at every symmetry-equivalent position on the particle. Prior
work on hallucinating cyclic symmetric oligomers
\parencite{wickyHallucinatingSymmetricProtein2022}, scaffolding at the scale of the
structural universe \parencite{lin_out_2024}, and integrating functional motifs
into designed nanoparticles \parencite{haas_novo_2026} indicates this is within
reach once the whole-assembly modelling barrier is removed. This work will be
carried out in collaboration with the Baker Lab during a visit there in the next
academic year (2026/27), giving direct access to their design and experimental
pipelines for validating conditionally designed immunogens.

\section{Greater control over the shape of generated assemblies}
\label{sec:shape-control}

A second direction is to gain finer control over the \emph{shape} of the generated
proteins and assemblies (their overall size, envelope and geometry) independently
of, or jointly with, the motif constraint. Existing generative models already offer
some spatial and shape conditioning: Chroma exposes programmable shape and
symmetry constraints \parencite{ingrahamIlluminatingProteinSpace2023}, ProtComposer
conditions structure generation on 3D ellipsoids
\parencite{stark_protcomposer_2025}, and ProxelGen generates proteins as 3D
densities \parencite{faltings_proxelgen_2025}; related work steers a pretrained
protein diffusion model at inference time to match experimental measurements such
as cryo-EM density maps, either by guiding the atomic coordinates directly
\parencite{raghu_multiscale_2025} or by optimising the model's conditional
embeddings \parencite{li_robust_2026}. I plan to improve on this spatial/shape conditioning so
that a designer can specify the target envelope of a nanoparticle (for instance to
match a desired antigen valency or a delivery constraint) and have the model
respect it, ideally composed with the symmetric-scaffolding capability of
Section~\ref{sec:motif-scaffolding}.

\section{A broader agenda: designing large symmetric protein assemblies}
\label{sec:broad-agenda}

More broadly, and over the remainder of the DPhil, I will continue to study the
design of large, and often symmetric, protein assemblies as a subject in its own
right. Doing this well requires understanding, and learning to design for, the
factors that govern whether subunits actually come together: specific
protein--protein interactions at the designed interfaces and the
thermodynamics of self-assembly, the functional and physicochemical drivers of
symmetric oligomerisation \parencite{goodsell_structural_2000} and the multistep,
ordered assembly pathways documented for natural complexes
\parencite{marsh_structure_2015}, and phenomena such as the structural
flexibility and oligomorphism observed in designed assemblies
\parencite{khmelinskaia_local_2025}.
The space of targets is correspondingly wide, from vaccine scaffolds to
mechanically coupled molecular machines \parencite{courbet_computational_2022} and
allosterically switchable assemblies \parencite{pillai_novo_2024}. This is the
outermost ring of the funnel: a general research direction that the concrete
projects above both feed into and help define.

\section{Timeline and milestones}
\label{sec:gantt}

Figure~\ref{fig:gantt} summarises the plan on a single timeline: the work
packages of Sections~\ref{sec:motif-scaffolding}--\ref{sec:broad-agenda}, the
papers and conferences they are aimed at, and the two formal milestones of the
DPhil.

\clearpage % put the Gantt chart on its own page so it can't overlap the text above
\begin{figure}[H]
  \centering
  \hspace*{-1.2cm}% shift the chart left (tweak as needed)
  \begin{ganttchart}[
      hgrid, vgrid,
      x unit=1.45cm,
      y unit title=0.6cm,
      y unit chart=0.75cm,
      title height=1,
      title/.append style={fill=oxfordblue!12},
      title label font=\footnotesize\bfseries,
      bar/.append style={fill=oxfordblue!75, rounded corners=1.5pt},
      bar height=0.55,
      bar label font=\footnotesize,
      bar label node/.append style={text width=5.2cm, align=right},
      group/.append style={fill=oxfordblue},
      group height=0.55,
      group label font=\footnotesize\bfseries,
      group label node/.append style={text width=5.2cm, align=right},
      milestone/.append style={fill=orange!85},
      milestone height=0.55,
      milestone label font=\footnotesize\itshape,
      milestone label node/.append style={text width=5.2cm, align=right},
      group right shift=0,
      group left shift=0,
    ]{1}{8}
    % --- title rows: years and quarters (TODO: adjust to your real timeline) ---
    \gantttitle{2026}{2}\gantttitle{2027}{4}\gantttitle{2028}{2} \\
    \gantttitle{Q3}{1}\gantttitle{Q4}{1}%
    \gantttitle{Q1}{1}\gantttitle{Q2}{1}\gantttitle{Q3}{1}\gantttitle{Q4}{1}%
    \gantttitle{Q1}{1}\gantttitle{Q2}{1} \\
    % --- Foundation: whole-assembly modelling ---
    \ganttgroup{Whole-assembly modelling}{1}{5} \\
    \ganttbar{Finalise Design-CP \& submit paper}{1}{2} \\
    \ganttbar{Experimental (in-vitro) validation}{3}{5} \\
    % --- Conditional generation ---
    \ganttgroup{Conditional generation}{3}{8} \\
    \ganttbar{Symmetric motif scaffolding\\(Baker Lab visit)}{3}{6} \\
    \ganttbar{Test on real antigens\\(Baker Lab visit)}{6}{8} \\
    % --- Shape / spatial control ---
    \ganttgroup{Shape \& spatial control}{1}{4} \\
    \ganttbar{Improve spatial/shape\\conditioning}{1}{4} \\
    % --- Broader agenda ---
    \ganttgroup{Large symmetric assemblies}{6}{8} \\
    \ganttbar{PPIs \& thermodynamics\\of assembly}{6}{8} \\
    \ganttbar{Thesis writing}{7}{8} \\
    % --- Papers and conferences ---
    % Each paper bar ends on the submission deadline of the conference it targets:
    % ICML 2027 closes in 2027 Q1, ICLR 2028 in 2027 Q3.
    \ganttgroup{Papers \& conferences}{2}{8} \\
    \ganttbar{Shape-conditioning paper\\(ICML 2027)}{2}{3} \\
    \ganttbar{Symmetric-scaffolding paper\\(ICLR 2028)}{4}{5} \\
    \ganttmilestone{ICML 2027}{4} \\
    \ganttmilestone{ICLR 2028}{7} \\
    % --- Milestones ---
    \ganttmilestone{Transfer of Status}{1} \\
    \ganttmilestone{Confirmation of Status}{5}
  \end{ganttchart}
  \caption{Gantt chart of upcoming milestones and deadlines}
  \label{fig:gantt}
\end{figure}
